\deleted{We proposed an edge-computing system that estimates the activity levels of people working in office environments using a LiDAR-sensor network.  



The estimation results demonstrated that three-dimensional features derived from point-cloud data are effective for activity-level estimation.

  %

  The XGBoost model achieved its highest accuracy when all features were used as inputs.

% 

To validate the system, experiments were conducted in an office environment with multiple participants. 

% 

The collected point-cloud data were transformed into statistical features that represented different aspects of participants' movements. 

% 

To estimate activity levels, we constructed two types of machine learning models that used either single point-cloud features or all features.  

For the XGBoost model, incorporating all features achieved higher accuracy than any single-feature model.

%

For the RF model, image-based 50th percentile values achieved the highest accuracy.

%

We considered the impact of clustering instability on the estimation accuracy.}

\revised{This paper proposed the system that estimates activity levels of office workers using a LiDAR-sensor network.

Three-dimensional features derived from point-cloud data were effective for activity-level estimation.

The XGBoost model achieved the highest accuracy with all features as inputs.

The RF model achieved the highest accuracy with image-based 50th percentile values.

Clustering instability was identified as a factor limiting estimation accuracy improvement.}

\revised{Office experiments collected point-cloud data from 37 participants during SQL learning-video sessions.

Point-cloud data were converted into image-based and centroid-based statistical features.

RF and XGBoost classifiers predicted ML agreement from feature-vector differences between session pairs.}

% \begin{revisedblock}

% %1. 評価の限界

% The evaluation was limited to SQL learning-video sessions with student participants.

%   %1-1. 適用できる範囲

%   The estimation accuracy applies to seated PC-based deskwork of the evaluated type.

% %2. 未評価

% Other seated deskwork in operational offices was not evaluated.

%   %2-1. 今後

%   Future evaluation will examine other seated deskwork in operational offices.

% \end{revisedblock}

